\chapter*{References and Bibliography}
\addcontentsline{toc}{chapter}{References and Bibliography}

This chapter outlines the foundational research, software libraries, and theoretical frameworks that have guided the development of the Hybrid Spatial-Frequency Image Classification project. The references encompass the core deep learning architectures, classical signal processing algorithms, explainability mechanisms, and the datasets utilized for benchmarking.

\renewcommand{\bibname}{} % Removes the default "Bibliography" title so it doesn't clash with our custom chapter title above
\vspace{-2em} % Reduces the gap left by the removed default title

\begin{thebibliography}{99}

% ==========================================
% DEEP LEARNING ARCHITECTURES & ALGORITHMS
% ==========================================

\bibitem{resnet}
He, K., Zhang, X., Ren, S., \& Sun, J. (2016). 
\textit{Deep Residual Learning for Image Recognition}. 
In Proceedings of the IEEE conference on computer vision and pattern recognition (pp. 770-778).

\bibitem{xgboost}
Chen, T., \& Guestrin, C. (2016). 
\textit{XGBoost: A Scalable Tree Boosting System}. 
In Proceedings of the 22nd ACM SIGKDD International Conference on Knowledge Discovery and Data Mining (pp. 785-794).

\bibitem{freq_learning}
Stuchi, J. A., Boccato, L., \& Attux, R. (2020). 
\textit{Frequency learning for image classification}. 
arXiv preprint arXiv:2006.15476.

% ==========================================
% SIGNAL PROCESSING & MATHEMATICAL FOUNDATIONS
% ==========================================

\bibitem{fft}
Cooley, J. W., \& Tukey, J. W. (1965). 
\textit{An algorithm for the machine calculation of complex Fourier series}. 
Mathematics of computation, 19(90), 297-301.

\bibitem{pca}
Wold, S., Esbensen, K., \& Geladi, P. (1987). 
\textit{Principal component analysis}. 
Chemometrics and intelligent laboratory systems, 2(1-3), 37-52.

\bibitem{gonzalez_woods}
Gonzalez, R. C., \& Woods, R. E. (2018). 
\textit{Digital Image Processing} (4th ed.). 
Pearson.

% ==========================================
% EXPLAINABILITY & MODEL INTERPRETATION
% ==========================================

\bibitem{gradcam}
Selvaraju, R. R., Cogswell, M., Das, A., Vedantam, R., Parikh, D., \& Batra, D. (2017). 
\textit{Grad-CAM: Visual Explanations from Deep Networks via Gradient-based Localization}. 
In Proceedings of the IEEE international conference on computer vision (pp. 618-626).

\bibitem{gradcam_plus_plus}
Chattopadhay, A., Sarkar, A., Howlader, P., \& Balasubramanian, V. N. (2018). 
\textit{Grad-CAM++: Generalized gradient-based visual explanations for deep convolutional networks}. 
In 2018 IEEE winter conference on applications of computer vision (WACV) (pp. 839-847).

\bibitem{scorecam}
Wang, H., Wang, Z., Du, M., Yang, F., Zhang, Z., Ding, S., ... \& Hu, X. (2020). 
\textit{Score-CAM: Score-weighted visual explanations for convolutional neural networks}. 
In Proceedings of the IEEE/CVF conference on computer vision and pattern recognition workshops (pp. 24-25).

\bibitem{cam}
Zhou, B., Khosla, A., Lapedriza, A., Oliva, A., \& Torralba, A. (2016). 
\textit{Learning Deep Features for Discriminative Localization}. 
In Proceedings of the IEEE conference on computer vision and pattern recognition (pp. 2921-2929).

\bibitem{shap}
Lundberg, S. M., \& Lee, S. I. (2017). 
\textit{A unified approach to interpreting model predictions}. 
In Advances in neural information processing systems (pp. 4765-4774).

\bibitem{lime}
Ribeiro, M. T., Singh, S., \& Guestrin, C. (2016). 
\textit{"Why should I trust you?" Explaining the predictions of any classifier}. 
In Proceedings of the 22nd ACM SIGKDD international conference on knowledge discovery and data mining (pp. 1135-1144).

% ==========================================
% DATASETS & COMPUTER VISION BENCHMARKS
% ==========================================

\bibitem{stl10}
Coates, A., Ng, A., \& Lee, H. (2011). 
\textit{An Analysis of Single-Layer Networks in Unsupervised Feature Learning}. 
In Proceedings of the fourteenth international conference on artificial intelligence and statistics (pp. 215-223). JMLR Workshop and Conference Proceedings.

\bibitem{vggface}
Parkhi, O. M., Vedaldi, A., \& Zisserman, A. (2015). 
\textit{Deep Face Recognition}. 
In British Machine Vision Conference.

% ==========================================
% SOFTWARE LIBRARIES & FRAMEWORKS
% ==========================================

\bibitem{pytorch}
Paszke, A., Gross, S., Massa, F., Lerer, A., Bradbury, J., Chanan, G., ... \& Chintala, S. (2019). 
\textit{PyTorch: An Imperative Style, High-Performance Deep Learning Library}. 
In Advances in Neural Information Processing Systems (pp. 8024-8035).

\bibitem{scikitlearn}
Pedregosa, F., Varoquaux, G., Gramfort, A., Michel, V., Thirion, B., Grisel, O., ... \& Duchesnay, E. (2011). 
\textit{Scikit-learn: Machine learning in Python}. 
Journal of machine learning research, 12(Oct), 2825-2830.

\bibitem{opencv}
Bradski, G. (2000). 
\textit{The OpenCV Library}. 
Dr. Dobb's Journal of Software Tools.

\bibitem{albumentations}
Buslaev, A., Iglovikov, V. I., Khvedchenya, E., Parinov, A., Druzhinin, M., \& Kalinin, A. A. (2020). 
\textit{Albumentations: fast and flexible image augmentations}. 
Information, 11(2), 125.

\end{thebibliography}
